% -------------------------------------------------------------------------
% WBS ID: RES-GEO-PAR
% Deliverable: Parametric Geometry and Design-Variable Definition
% Status: In progress
% Evidence status: Active comparison descriptor set established
% Review status: Not reviewed
% Last updated: 2026-07-19
% Dependencies: RES-GEO-TAX, RES-DAT-SCN, RES-NUM-MSH
% Downstream interfaces: RES-GEO-MET, RES-GEO-REP, Software geometry manifest
% -------------------------------------------------------------------------

\section{RES-GEO-PAR --- Parametric Geometry and Design-Variable Definition}
\label{sec:res-geo-par}

\subsection{Purpose and Scope}
\label{sec:res-geo-par-purpose}

This Deliverable defines the bounded parameters required to compare
fixed-rigid barrier classes in the two local corridors. It does not
define an optimisation-ready free-form design space and does not
activate ML or generative geometry.

\subsection{Research Questions}
\label{sec:res-geo-par-research-questions}

\begin{enumerate}
  \item Which geometric parameters are required for wall-type and
    obstacle/array-type barriers?
  \item Which parameters are measured, inferred, scenario-controlled or
    unresolved?
  \item Which parameter bounds are controlled by terrain coverage,
    corridor width and meshability rather than design preference?
\end{enumerate}

\subsection{Terminology and Definitions}
\label{sec:res-geo-par-definitions}

% TODO: Define the technical terminology, symbols, quantities and
% classifications used in this Deliverable.

\subsection{Literature Evidence}
\label{sec:res-geo-par-literature-evidence}

\textcite{deljesus2012porous} supports separate parameter treatment
for impermeable walls and barriers whose openings or porosity change
hydraulic response. Source role: supports explicit opening/spacing
parameters and limits unresolved porous claims. \textcite{watanabe2022validation}
supports simple controlled geometry parameters for local impact and
loading validation. Project conclusion: the parameter set shall remain
small, inspectable and tied to hydrodynamic metrics.

\subsection{Governing Relations and Quantitative Definitions}
\label{sec:res-geo-par-governing-relations}

For each case $j$, the comparison parameter vector is:

\begin{align}
  \boldsymbol{\theta}^{\mathrm{geo}}_j
  &=
  \left[
    H_j,\,
    C_j,\,
    B_j,\,
    L_j,\,
    \psi_j,\,
    x'_j,\,
    y'_j,\,
    S_j,\,
    N_j,\,
    \phi^{\mathrm{open}}_j
  \right]^{\mathsf{T}}
  \label{eq:res-geo-par-parameter-vector}
\end{align}

where $\phi^{\mathrm{open}}_j$ is the resolved open-area or gap
fraction for array or segmented cases. Source role:
\textcite{deljesus2012porous} supports separating impermeable and
opening-controlled barriers; \textcite{watanabe2022validation}
supports maintaining geometry parameters that map directly to local
loading observables. Limitation: \cref{eq:res-geo-par-parameter-vector}
does not set final bounds.

\subsection{Geometry Family or Representation}
\label{sec:res-geo-par-geometry-family-representation}

% TODO: Complete this RES-GEO Domain-specific subsection.

\subsection{Design Variables and Parameter Bounds}
\label{sec:res-geo-par-design-variables-parameter-bounds}

Parameter sets shall distinguish measured geometry, inferred geometry,
scenario-controlled geometry and deferred structural assumptions.
Candidate parameters may include crest height, barrier height, wall or
element thickness, footprint length, plan orientation, corridor
placement, gap spacing, row count, resolved opening fraction and
surface roughness flag. Foundation strengthening, material changes,
damage evolution and deformable response are stretch extensions owned
by RES-STR after hydrodynamic load extraction is stable.

\subsection{Geometric Descriptors}
\label{sec:res-geo-par-geometric-descriptors}

% TODO: Complete this RES-GEO Domain-specific subsection.

\subsection{Placement and Arrangement Rules}
\label{sec:res-geo-par-placement-arrangement-rules}

% TODO: Complete this RES-GEO Domain-specific subsection.

\subsection{Feasibility Constraints}
\label{sec:res-geo-par-feasibility-constraints}

% TODO: Complete this RES-GEO Domain-specific subsection.

\subsection{Two- and Three-Dimensional Representation}
\label{sec:res-geo-par-two-three-dimensional-representation}

% TODO: Complete this RES-GEO Domain-specific subsection.

\subsection{Candidate Approaches}
\label{sec:res-geo-par-candidate-approaches}

% TODO: Define the credible candidate approaches considered.

\subsection{Critical Comparison}
\label{sec:res-geo-par-critical-comparison}

% TODO: Compare candidates using physically and project-relevant criteria.
% Do not select an approach solely on popularity or implementation ease.

\subsection{Assumptions and Applicability}
\label{sec:res-geo-par-assumptions}

% TODO: State assumptions and the conditions under which they are valid.

\subsection{Limitations and Failure Conditions}
\label{sec:res-geo-par-limitations}

% TODO: State where the evidence, model or selected approach ceases to be
% reliable or applicable.

\subsection{Project-Specific Conclusions}
\label{sec:res-geo-par-conclusions}

The active Studentship comparison requires a compact parameter vector,
not an open-ended design grammar. Geometry choices must be reproducible
from the manifest, meshable in the local domain and traceable to
barrier-performance metrics.

\subsection{Selected Approach and Rationale}
\label{sec:res-geo-par-selected-approach}

The selected approach is a descriptor-first parameter set for fixed
wall, obstacle, dissipating and array barriers. Broad geometric
optimisation and generative parameter discovery are formally deferred.

\subsection{Unresolved Decisions}
\label{sec:res-geo-par-unresolved-decisions}

\begin{itemize}
  \item Final numeric bounds for $H_j$, $B_j$, $L_j$, $S_j$ and
    $N_j$ remain open until corridor terrain and mesh constraints are
    accepted.
  \item Real-defence reconstruction parameters for Kamaishi remain
    TODO pending a citable geometry source.
  \item Roughness and opening/porosity flags remain provisional until
    RES-MOD-SRC decides whether any porous closure is active.
\end{itemize}

\subsection{Requirements and Handoffs}
\label{sec:res-geo-par-requirements}

Software shall ingest \cref{eq:res-geo-par-parameter-vector} from the
geometry manifest, generate fixed barrier patches, and preserve stable
case identifiers for metric extraction. RES-NUM-MSH shall report
whether each parameter set is meshable; RES-VER-MET shall record which
geometry parameters produced each run-up, overtopping, force, moment
and energy metric.

\subsection{Deliverable Completion Record}
\label{sec:res-geo-par-completion-record}

% TODO: Record whether the Deliverable is:
%
% - Not started
% - In progress
% - Draft complete
% - Reviewed
% - Accepted
%
% Acceptance requires:
%
% - the research questions to be answered;
% - claims to be supported by appropriate evidence;
% - assumptions and limitations to be explicit;
% - the project-specific conclusion to be stated;
% - downstream requirements to be traceable;
% - unresolved decisions to be closed or formally deferred.
